% Source: https://artofproblemsolving.com/community/c3h3039567_simple_bijection
% Stable identity: aops-topic-3039567-q01
(a) The divisor involution $d\mapsto \dfrac{n^2}{d}$ pairs divisors below $n$ with those above it, leaving only $n$ fixed. There are $(28+1)(34+1)(6+1)(4+1)=35525$ divisors, so the answer is $\boxed{17762}$. (b) Similarly $n^4$ has $(56+1)(68+1)(12+1)(8+1)=460161$ divisors, with sole fixed divisor $n^2$, giving $\boxed{230080}$. (c) The cutoff $n^2$ is not the square root of $n^3$, so the same halving argument does not apply. Every divisor is uniquely $2^a3^b5^c7^d$ with $0\le a\le42,0\le b\le51,0\le c\le9,0\le d\le6$. Exact enumeration of this finite domain gives $\boxed{126480}$ below $n^2$.
\begin{verbatim}
n=2**14*3**17*5**3*7**2
count=sum(2**a*3**b*5**c*7**d<n*n
          for a in range(43) for b in range(52)
          for c in range(10) for d in range(7))
assert count==126480
\end{verbatim}
